\section{Drupal Project Analysis} % {{{

\marginnote{\includegraphics[width=\marginparwidth]{drupal/logo}}

Drupal is a \ac{CMS} and a \ac{CMF}. It is written in the PHP programming
language and licensed under the \ac{GPL}. Drupal is used for a wide range
of capabilities such as blogs, company and political websites, communities and
is directed to be a widely spread cms \cite{DrupalOverview}. Drupal is also
known as \emph{Drupal Core} inside the Community and contains the base
functionality of the cms. However it can be extended through numerous modules.
In the following though only Drupal Core will be analysed.

\begin{figure}[htbp]
  \centering
  \includegraphics[width=0.8\textwidth]{drupal/commits_by_author}
  \caption{Monthly activity of the most active Drupal Core developers}
\end{figure}

\subsection{History} % {{{

The first version of Drupal was released in 2001 by Dries Buytaert
\cite{DrupalHistory}. Originally the software took action in providing a black
board like directory. However quickly a fully-fledged \ac{CMS} emerged out of
the initial version. The name originates from the dutch word \emph{druppel}
which translated has the meaning of a drop. Especially in the last years it
experienced a wide distribution and is used for at least \unit[1.7]{\%} of all
websites worldwide and holds a marketshare of \unit[6.2]{\%}
\cite{DrupalBuiltWith,DrupalW3Techs}.

% }}}

\subsection{Community} % {{{

The Drupal project has a large user and developer community available.
According to own statements, over 250.000 users were registered on
\url{drupal.org} in August 2011. About 1200 of those were additionally listed
as developers, whereas the number might be a bit higher \cite{DrupalBuytaert}.

Twice a year the community holds a conference under the name \emph{Drupal
Conference} or \emph{DrupalCon}. To alleviate the journey of the attendants the
conference takes place alternatively in Europe and North America.

The communication possibilities sprawl from forums on \url{drupal.org}, mailing
lists and discussion groups on groups.drupal.org. Additionally there exist a
wide range of \ac{IRC} channels for different aspects of the Drupal project.

\begin{table}
  \centering
  \begin{tabularx}{\textwidth}{XXl}
    \toprule
    \tableheadline{Venue} & \tableheadline{Date}  & \tableheadline{Attendees} \\
    \midrule
    Munich                & August 2012           & N/A \\
    Denver                & March 2012            & N/A \\
    London                & August 2011           & 1751 \\
    Chicago               & March 2011            & 3000 \\
    Kopenhagen            & August 2010           & 1200 \\
    San Francisco         & April 2010            & 3000 \\
    Paris                 & September 2009        & 850 \\
    Washington, D.C.      & March 2009            & 1400 \\
    Szeged                & August 2008           & 500 \\
    Boston                & March 2008            & 850 \\
    Barcelona             & September 2007        & 450 \\
    Sunnyvale             & March 2007            & \textasciitilde 300 \\
    Brussels              & September 2006        & 150 \\
    Vancouver             & February 2006         & \textasciitilde 150 \\
    Amsterdam             & October 2005          & \textasciitilde 100 \\
    Portland              & August 2005           & 100 \\
    Antwerpen             & February 2005         & \textasciitilde 50 \\
    \bottomrule
  \end{tabularx}
  \caption{Previous and planned DrupalCon events \cite{DrupalWalling}}
\end{table}

Next to the last-mentioned classification of users and developers, people can
-- provided that they are working on Drupal Core -- classified into one of the
following categories \cite{DrupalCoreDevelopers}.

\begin{figure}[htbp]
  \centering
  \includegraphics[width=0.8\textwidth]{drupal/commits_by_year}
  \caption{Yearly overview of commits to Drupal Core}
\end{figure}

\begin{description}

  \item[Founder and Lead Developer] Dries Buytaert, the founder of the Drupal
    project, holds the role of the project leader and therefore decides about
    the direction in which the project is heading to. Furthermore he is the
    one, who most often takes a decision and also often makes the last
    decision. Though in some cases he also gives away this control to a trusted
    person.

  \item[Core Committer] Only few people have write access to the Drupal Core
    repository. They are known as Core Committer and look through code changes
    and propositions and decide about acceptance of a patch. At this juncture
    there is a distionction between permanent Core Committers and Branch
    Maintainers.

    \begin{description}
      \item[Permanent Core Committer] Dries Buytaert
      \item[Branch Maintainer] Next to Dries Buytaert they are in charge of
      maintaining a particular Drupal version.
      \begin{itemize}
        \item Gerhard Killesreiter for Drupal 4.7.x.
        \item Neil Drumm for Drupal 5.x.
        \item Gábor Hojtsy for Drupal 6.x.
        \item Angela Byron for Drupal 7.x.
      \end{itemize}
    \end{description}

  \item[Maintainer] While they not have to be involved in the decision making
    progress, maintainer have the liability of small portions of Drupal Core.
    Most of them are particular Core modules or technical domains such as
    JavaScript. The assignment happens over Dries Buytaert, whom interested
    developers can approach or will be invited to join.

  \item[Core Contributor] All developers who bring in code or documentation for
    Drupal Core are known as Core Contributors. All modification proposals are
    looked through another Core Committer in a review process and then either
    accepted or dropped.

\end{description}

% }}}

\subsection{Release Process} % {{{

For releases, Drupal uses a version naming with two numbers since Drupal 5. The
numbers stand for major and minor versions \cite{DrupalUpgrade}. New major
releases are quite scarce and will get released every few years. They are
planned for a long time in advance, often with incompatible changes. Before
Drupal 5 a three numbers versioning scheme was used.

While the time periods between single Drupal Core major releases are not bound,
the time until a stable release is divided into single phases
\cite{DrupalReleaseCycle}. The length of a single phase is additionally not
fixated either but limited by a range of factors which will be explained in the
following.

\begin{figure}[htbp]
  \centering
  \includegraphics[width=\textwidth]{drupal/releases}
  \caption{Major Releases of Drupal}
\end{figure}

\begin{figure}[htbp]
  \centering
  \includegraphics[width=\textwidth]{drupal/punchcard}
  \caption{Time based view on commits of Core Contributors}
\end{figure}

\begin{description}

  \item[Code Thaw] After each major release a new branch will be created in the
    Drupal Core repository. In this phase all new features and modifications
    can be discussed and added. No restrictions apply.

  \item[Code Slush] This phase was added to the Drupal 7 release cycle to
    restrict the time between major releases. Important features will be
    elaborated and with the help of so called initiatives handled. In this
    phase no new features except the chosen initiatives are allowed. Though
    changes will be accepted if the fix errors or add missing functionality.

    At the end of this phase, the \ac{API} freeze will be declared. At this
    point in time, the \ac{API} can not be changed anymore except critical
    errors are found and have to be corrected.

\begin{figure}[htbp]
  \centering
  \includegraphics[width=0.8\textwidth]{drupal/release_phases}
  \caption{Drupal Release Phases}
\end{figure}

  \item[Polish Phase] In this phase Drupal Core will be polished and the last
    enhancements on performance, accessibility, documentation and other are
    brought in. At the end of the phase string freeze and user interface freeze
    will be in place.

  \item[Release Phase] Lastly in this phase there will be several alpha, beta,
    release candidates and in the long run the final release released. The
    first release candidate will be published once there are no more critical
    bugs present. Additional release candidates will be released if new
    critical bugs appear, otherwise, if no new critical bugs appear, the final
    release is published.

\end{description}

\begin{figure}[htbp]
  \centering
  \includegraphics[width=0.8\textwidth]{drupal/commits_by_month}
  \caption{Amount of commits per month of Core Contributors}
\end{figure}

% }}}

\subsection{Development} % {{{

The Drupal project changed it's development workflow for the upcoming Drupal 8
release. Until then, the development was mostly driven through the issue
queues, which is a bug tracking system on \url{drupal.org}. While it is still
used for communication, bugs and new features, the development process features
a number of core initiatives which will stand for a major area of Drupal
\cite{DrupalInitiatives}. Each initiative has one or two initiative owner who
leads the development and coordinates the initiative while working closely with
Dries Buytaert.

Initiatives can be created by any Drupal developer, however Dries Buytaert
choses which initiative will be followed for the next major release of Drupal.
He also choses the initiative leader and remains in close contact with him.
Furthermore each initiative has defined a small set of goals which should be
reached by the time of the next major release. For the Drupal 8 release, the
following initiatives will be followed.

\begin{itemize}
  \item Configuration Management
  \item Web Services Context Core Initiative
  \item Design 4 Drupal
  \item Drupal 8 Multi-Lingual Initiative
  \item HTML5
  \item Mobile
\end{itemize}

\begin{figure}[htbp]
  \centering
  \includegraphics[width=0.8\textwidth]{drupal/authors_by_month}
  \caption{Amount of distinct authors of Drupal Core over time}
\end{figure}

% }}}

% }}}
